% =============================================================================
%  Notas de clase — Sesión 17: Taller 1 — Resolución de problemas
%  Programación Avanzada — Universidad Panamericana
%  Compilar: pdflatex sesion17_taller_1.tex
% =============================================================================
\documentclass[12pt, oneside]{article}
\usepackage{progavanzada}
\usepackage{array}

\begin{document}

\portada{Sesión 17}{Taller 1 — Resolución de problemas}{2026}

\tableofcontents
\newpage

% =============================================================================
\section{Los árbitros automáticos}
% =============================================================================

Antes de que existieran los jueces en línea, los concursos de programación se
calificaban a mano: un humano ejecutaba tu código, leía la salida y decidía si
estaba bien.  En 1989, la Universidad de Valladolid montó el primer
\textbf{juez automático} (\textit{Online Judge}) para sus prácticas de
algoritmia.  La idea era simple: el juez tiene los casos de prueba correctos;
tú mandas tu código; él lo compila, lo ejecuta con esos casos y compara la
salida carácter por carácter.

Hoy existen cientos de jueces: Codeforces, LeetCode, HackerRank, SPOJ,
y \textbf{omegaUp}, la plataforma mexicana usada en la Olimpiada Mexicana de
Informática.  El veredicto que ves después de enviar puede ser:

\begin{center}
\begin{tabular}{ll}
  \toprule
  \textbf{Veredicto} & \textbf{Significado} \\
  \midrule
  \texttt{AC} — Accepted            & Tu salida es exactamente la esperada \\
  \texttt{WA} — Wrong Answer        & Tu salida difiere en al menos un carácter \\
  \texttt{TLE} — Time Limit Exceeded & Tu programa tardó más del límite permitido \\
  \texttt{RE} — Runtime Error       & Tu programa terminó con error (e.g.\ segfault) \\
  \texttt{CE} — Compilation Error   & El código no compiló \\
  \bottomrule
\end{tabular}
\end{center}

\begin{advertencia}
  El juez compara tu salida \textbf{carácter por carácter}.  Un espacio
  de más o un salto de línea faltante puede darte \texttt{WA} aunque la
  lógica sea correcta.  Lee el formato de salida con cuidado.
\end{advertencia}

% =============================================================================
\section{Instrucciones del taller}
% =============================================================================

Esta sesión es trabajo en \textbf{equipos de dos o tres personas}.
Cada problema tiene dos partes:

\begin{description}
  \item[Parte 1 — Revisión guiada.]
    El profesor resuelve el caso de prueba del enunciado en el pizarrón.
    Tu tarea en esta parte es \textbf{poner atención}: observa cómo se
    lee la entrada línea por línea, qué operaciones se hacen con esos
    datos y cómo se obtiene exactamente la salida esperada.
    No abras el cuaderno todavía.

  \item[Parte 2 — Trabajo en equipo.]
    Con lo observado en la Parte 1, tu equipo encuentra la condición
    matemática general, escribe el pseudocódigo, implementa en C++ y
    envía al juez.
\end{description}

\begin{advertencia}
  Durante la \textbf{Parte 1} no escribas código ni intentes adelantarte
  a la solución.  El objetivo es entender \emph{qué} pide el problema,
  no todavía \emph{cómo} resolverlo para cualquier entrada.
\end{advertencia}

\begin{consejo}
  \textbf{Flujo de la Parte 2 para cada problema:}
  \begin{enumerate}
    \item Con el ejemplo ya resuelto, generaliza: ¿qué condición aplica
          para cualquier entrada?
    \item Escribe el algoritmo en pseudocódigo (una oración por paso).
    \item Implementa en C++ y prueba con el mismo ejemplo del enunciado.
    \item Envía al juez.  Si recibes \texttt{WA}, construye un
          contraejemplo antes de cambiar código al azar.
  \end{enumerate}
\end{consejo}

\newpage

% =============================================================================
\section{Problema 1: Paradoja de los promedios}
% =============================================================================

\subsection{Enunciado}

Un chiste conocido en Ciencias de la Computación dice: si un mal estudiante de
Cómputo se cambia a la facultad de Economía, \textit{sube el promedio de
inteligencia en ambas facultades al mismo tiempo}.

Dado el IQ de todos los estudiantes de las dos facultades, determina cuántos
estudiantes de Cómputo podrían protagonizar el chiste.

\textbf{Entrada.}
La primera línea contiene un entero $T$ — el número de casos de prueba.
Cada caso de prueba es precedido por una línea en blanco.

Cada caso comienza con dos enteros positivos $n$ y $m$ — el número de
estudiantes de Cómputo y de Economía, respectivamente.
El número de estudiantes de Cómputo es al menos 2.

Las siguientes líneas contienen $n + m$ enteros positivos separados por
espacios: los primeros $n$ son los IQs de los estudiantes de Cómputo;
los $m$ restantes corresponden a Economía.

\textbf{Salida.}
Por cada caso de prueba, imprime en una sola línea el número de estudiantes
de Cómputo que harían que el chiste sea verdad.

\subsection{Parte 1 — Revisión guiada}

\begin{advertencia}
  \textbf{Pon atención al pizarrón.}  El profesor va a leer la entrada
  línea por línea y a mostrar, paso a paso, cómo obtener la salida
  \texttt{1} para el siguiente caso de prueba.  Observa qué operaciones
  se hacen con los números antes de poder contestar.
\end{advertencia}

\begin{center}
\begin{tabular}{ll}
  \toprule
  \textbf{Entrada} & \textbf{Salida} \\
  \midrule
  \begin{minipage}[t]{4cm}
    \texttt{1}\\
    \texttt{}\\
    \texttt{5 5}\\
    \texttt{100 101 102 103 104}\\
    \texttt{98 100 102 99 101}
  \end{minipage}
  &
  \begin{minipage}[t]{4cm}
    \texttt{1}
  \end{minipage} \\
  \bottomrule
\end{tabular}
\end{center}

\textbf{Notas de lo observado en el pizarrón:}

\vspace{3cm}

\subsection{Parte 2 — Trabajo en equipo}

\subsubsection{Análisis}

¿Qué condición debe cumplir un estudiante de Cómputo con IQ $x$ para que
al cambiar de facultad \textbf{suba el promedio de Cómputo}?

\vspace{3cm}

¿Qué condición debe cumplir el mismo estudiante para que al llegar a Economía
\textbf{suba el promedio de Economía}?

\vspace{3cm}

Si el promedio de Cómputo es $\bar{c}$ y el de Economía es $\bar{e}$,
¿en qué rango debe estar $x$ para cumplir \textit{ambas} condiciones
simultáneamente?

\vspace{2cm}

\begin{recuerda}
  Comparar con punto flotante puede introducir errores de redondeo.  Si
  $\bar{c} = \text{sum}(c)/n$, la condición $x < \bar{c}$ es equivalente a
  $x \cdot n < \text{sum}(c)$ — todo con aritmética entera.
\end{recuerda}

\subsubsection{Pseudocódigo}

\vspace{5cm}

\newpage

% =============================================================================
\section{Problema 2: Ayudando a mi hermanito}
% =============================================================================

\subsection{Enunciado}

\textbf{(omegaUp: \texttt{La-tarea-de-mi-hermanito})}

Mi hermano está aprendiendo a ubicar números de tres dígitos en la recta
numérica.  Para ayudarle a practicar, le propongo un juego: le doy dos
números; él los tiene que leer al revés (de derecha a izquierda) y debe
decir cuál de los dos, leídos en reversa, es mayor.

Escribe un programa que revise las respuestas de tu hermano.

\textbf{Entrada.}
Dos números de tres dígitos separados por un espacio, $A$ y $B$.
$A$ y $B$ no pueden ser iguales y no contienen ceros.

\textbf{Salida.}
El número más grande de los dos \textit{leídos en reversa}, escrito en
reversa — exactamente como tu hermano lo debe leer.

\subsection{Parte 1 — Revisión guiada}

\begin{advertencia}
  \textbf{Pon atención al pizarrón.}  El profesor va a mostrar cómo
  leer los dos números, qué operación se aplica a cada uno y cómo se
  decide qué imprimir.  Observa en qué orden llegan los datos y qué
  tiene que aparecer exactamente en la salida.
\end{advertencia}

\begin{center}
\begin{tabular}{ll}
  \toprule
  \textbf{Entrada} & \textbf{Salida} \\
  \midrule
  \texttt{735 892} & \texttt{537} \\
  \texttt{331 341} & \texttt{143} \\
  \bottomrule
\end{tabular}
\end{center}

\textbf{Notas de lo observado en el pizarrón:}

\vspace{3cm}

\subsection{Parte 2 — Trabajo en equipo}

\subsubsection{Análisis}

Verifica a mano los dos ejemplos.  Para cada número de tres dígitos $d_1 d_2 d_3$,
¿cuál es su reversa en términos de sus dígitos?

\vspace{2cm}

¿Cuál es la fórmula matemática para obtener la reversa de un número $n$ de
tres dígitos a partir de sus dígitos individuales?
(\textit{Pista: usa división entera y módulo.})

\vspace{3cm}

¿Hay algún caso especial que considerar dado que los números \textbf{no
contienen ceros}?

\vspace{2cm}

\subsubsection{Pseudocódigo}

\vspace{5cm}

\newpage

% =============================================================================
\section{Resumen de la sesión}
% =============================================================================

\begin{itemize}
  \item Un \textbf{juez en línea} califica tu código automáticamente comparando
        la salida carácter por carácter.
  \item El veredicto más importante es \textbf{AC} (tu solución es correcta);
        el más traicionero es \textbf{WA} (la lógica falla en algún caso).
  \item Para ``Paradoja de los promedios'': la condición es
        $\bar{e} < x < \bar{c}$, verificada con aritmética entera para
        evitar errores de punto flotante.
  \item Para ``Mi hermanito'': revertir un número de tres dígitos sin ceros
        es una operación de dígitos; no necesitas cadenas de texto.
\end{itemize}

\end{document}
